\chapter{Algorithms}

This appendix presents the key algorithms used in Ansieyes in pseudocode form. Due to the corporate nature of the project, actual source code has been omitted in favour of algorithmic descriptions.

\section{Issue Analysis Pipeline}

\begin{tcblisting}{listing only,colback=gray!10!white, breakable, boxsep=0pt,top=1mm,bottom=1mm,left=1mm,right=1mm,
listing options={
basicstyle=\small\ttfamily,
commentstyle=\color{gray},
backgroundcolor=\color{gray!25}}}
ALGORITHM: AnalyseIssue(title, description, codebase, retries)
--------------------------------------------------------------
INPUT:  Issue title, description, codebase text, max retries
OUTPUT: Structured analysis (type, severity, root cause, solutions)

1. Load codebase content from Repomix output file
2. Build prompt = SystemRole + title + description + codebase
                  + JSON format instructions
3. FOR attempt = 0 TO retries DO
     response = CallGeminiAPI(prompt)
     json_data = ExtractJSON(response)   -- try 3 strategies
     IF json_data is valid THEN
       analysis = ParseToModel(json_data)
     ELSE
       analysis = FallbackTextExtraction(response)
     END IF
     IF IsHighQuality(analysis) THEN
       RETURN analysis
     END IF
   END FOR
4. RETURN best available analysis
\end{tcblisting}

\section{Quality Validation}

\begin{tcblisting}{listing only,colback=gray!10!white, breakable, boxsep=0pt,top=1mm,bottom=1mm,left=1mm,right=1mm,
listing options={
basicstyle=\small\ttfamily,
commentstyle=\color{gray},
backgroundcolor=\color{gray!25}}}
ALGORITHM: IsHighQuality(analysis)
----------------------------------
INPUT:  Parsed analysis object
OUTPUT: Boolean (true if acceptable quality)

IF count(analysis.solutions) < 1       THEN RETURN false
IF analysis.confidence < 0.3           THEN RETURN false
IF analysis.summary.length < 20        THEN RETURN false
IF analysis.primary_cause.length < 10  THEN RETURN false
IF any file path contains "path/to/"   THEN RETURN false
IF any file path contains "example.py" THEN RETURN false
IF any solution description < 10 chars THEN RETURN false
RETURN true
\end{tcblisting}

\section{Prompt Injection Detection}

\begin{tcblisting}{listing only,colback=gray!10!white, breakable, boxsep=0pt,top=1mm,bottom=1mm,left=1mm,right=1mm,
listing options={
basicstyle=\small\ttfamily,
commentstyle=\color{gray},
backgroundcolor=\color{gray!25}}}
ALGORITHM: DetectInjection(text)
--------------------------------
INPUT:  Issue text to evaluate
OUTPUT: Risk level, confidence, detected patterns

-- Layer 1: ML-based
ml_score = PytectorModel.classify(text)
ml_detected = (ml_score > 0.5)

-- Layer 2: Pattern matching across 7 categories
patterns_found = []
FOR each category in [role, system, bypass, file,
                      code, extraction, leakage] DO
  FOR each regex in category.patterns DO
    IF regex matches text THEN
      patterns_found.add(category)
    END IF
  END FOR
END FOR
pattern_score = lookup_confidence(patterns_found)

-- Layer 3: Heuristic analysis
heuristic_score = 0
IF special_char_sequences(text) > 3    THEN score += 0.6
IF instruction_keywords(text) > 3      THEN score += 0.7
IF suspicious_delimiters(text) > 2     THEN score += 0.5
IF max_sentence_length(text) > 100     THEN score += 0.4

-- Aggregate
final_confidence = MAX(ml_score, pattern_score,
                       heuristic_score)
pattern_count = len(unique(patterns_found))
risk = ClassifyRisk(final_confidence, pattern_count)
RETURN InjectionResult(risk, final_confidence, patterns)
\end{tcblisting}

\section{Cosine Similarity Duplicate Detection}

\begin{tcblisting}{listing only,colback=gray!10!white, breakable, boxsep=0pt,top=1mm,bottom=1mm,left=1mm,right=1mm,
listing options={
basicstyle=\small\ttfamily,
commentstyle=\color{gray},
backgroundcolor=\color{gray!25}}}
ALGORITHM: CosineDuplicateDetect(new_title, new_desc,
                                 existing_issues, threshold)
-----------------------------------------------------
INPUT:  New issue title and description,
        list of existing issues, threshold (default 0.6)
OUTPUT: Duplicate detection result

1. FOR each issue (new + existing) DO
     text = lowercase(title + " " + title + " " + desc)
     text = remove_special_chars(text)
     text = normalise_whitespace(text)
   END FOR
2. Build TF-IDF matrix using unigrams and bigrams
   (max 5000 features, English stop words removed)
3. Compute cosine similarity between new issue
   vector and all existing issue vectors
4. Find max_similarity and best_match_index
5. IF max_similarity >= threshold THEN
     confidence = MIN(max_similarity * 1.2, 1.0)
     RETURN DuplicateResult(true, existing[best_match],
                            max_similarity, confidence)
   ELSE
     RETURN DuplicateResult(false, null,
                            max_similarity, max_similarity)
   END IF
\end{tcblisting}

\section{Librarian File Selection}

\begin{tcblisting}{listing only,colback=gray!10!white, breakable, boxsep=0pt,top=1mm,bottom=1mm,left=1mm,right=1mm,
listing options={
basicstyle=\small\ttfamily,
commentstyle=\color{gray},
backgroundcolor=\color{gray!25}}}
ALGORITHM: LibrarianFileSelection(title, desc, chunks)
------------------------------------------------------
INPUT:  Issue title, description, directory chunks
OUTPUT: List of relevant file paths

1. Load all .txt files from chunks directory
2. Send directory list to Gemini:
   "Which directories are relevant to this issue?"
3. Validate returned directory names against loaded chunks
4. FOR each relevant chunk DO
     Send chunk content to Gemini:
     "Which files in this directory are relevant?"
     FOR each returned path DO
       IF path has "/" AND has extension
          AND length < 200
          AND not a comment line THEN
         Add to relevant_files set
       END IF
     END FOR
   END FOR
5. RETURN relevant_files
\end{tcblisting}
